% 第 9 章。数字：checks/check_ch09.py；引文：data/snapshots/india_facts.csv（《汉书》《后汉书》的逐字核对见
% data/fetch/verify_quotes.py，其余见 data/fetch/verify_web_quotes.py）；印度出土的罗马钱：darley2014_coin_finds
% （先运行 data/fetch/darley2014_coin_finds.py）；图：figsrc/fig_ch09.py。
\chapter{古代印度与印度洋贸易}\label{ch:india}

\begin{chapterguide}
\guidequestion{罗马的金银为什么流向印度？季风贸易怎样把地中海、阿拉伯与南亚连成一个市场？}
\begin{guidegoals}
\item 描述打印币与《政事论》所反映的国家财政与市场管理；
\item 用季风贸易和罗马作者对贵金属外流的抱怨，说明古代的“国际收支”；
\item 比较印度、中国与罗马对贵金属的需求及其后果；
\item 评估这些证据在数量上的不确定性。
\end{guidegoals}
\guideprereq{\chref{ch:greekroman}\enspace\nameref{ch:greekroman}；\chref{ch:external}\enspace\nameref{ch:external}}
\end{chapterguide}

老普林尼（Pliny the Elder）在《自然史》里算过一笔账：“按最低的估计，印度、赛里斯（Seres）与阿拉伯半岛每年从我们的帝国拿走一亿塞斯特斯——我们为奢侈与女人付的代价就是这么高。”\cite{pliny}塞斯特斯（sestertius，拉丁文缩写 HS）是罗马的货币单位，4 塞斯特斯合 1 枚银币第纳尔（denarius）。中国的史书也记下了这门生意。《后汉书》说大秦（罗马）用金银作钱，“与安息、天竺交市于海中，利有十倍”\cite{houhanshu}。

一张 2 世纪中期的纸草，记下了一条船从印度运回的货：“赫尔玛波隆号（Hermapollon）船上运出的六批货，以银计”，共值 1154 塔兰特 2852 德拉克马\cite{casson1990}。塔兰特（talent）与德拉克马（drachma）是希腊的单位。塔兰特既是重量单位，也是钱的单位；这里是钱，1 塔兰特合 6000 德拉克马，这一船货合 692.7 万德拉克马。在安东尼时期的埃及，一枚四德拉克马（tetradrachm）与一枚第纳尔等值\cite{howgego2010}，所以 1 埃及德拉克马合 1 塞斯特斯。这样算，一条船的货就是普林尼那一亿的 6.9\%。

金银为什么流向印度？普林尼的数字可信吗？这一亿在罗马的经济里有多大？用今天的话说，罗马对东方是不是有一个“经常账户逆差”，即向外国买的多、卖给外国的少（\chref{ch:external}）？答案要从一年一个来回的季风航线说起。

\section{当时的世界}\label{sec:india-world}

罗马的钱运往印度，走的是海路。这些船从哪里出发，怎样横渡大洋，在印度的哪些港口靠岸？汉朝的使者走的又是哪一条路？

\subsubsection{一年一个来回}
伽卢斯（Gallus）任埃及长官时，斯特拉波（Strabo）和他在一起。斯特拉波在《地理学》里说，他那时了解到，“大约一百二十条船从米奥斯·霍尔莫斯（Myos Hormos）开往印度”，而在托勒密王朝的时候，几乎没有人敢走这条航线、做印度的生意\cite{strabo}。在另一卷里他又说：“从前，连二十条船都不敢驶入阿拉伯湾，或驶出湾口的海峡哪怕一点点；如今，大批船队被派往远至印度与埃塞俄比亚尽头的地方。”\cite{strabo}

普林尼把这条路一段一段写了下来（图~\ref{fig:india_route}(a)）。从亚历山大里亚附近沿尼罗河上行到科普托斯（Coptos），三百零八英里，季节风吹的时候十二天走完。从科普托斯穿过沙漠到红海边的贝勒尼基（Berenice），二百五十七英里，也要走十二天。旅客一般在仲夏、天狼星升起之前或刚升起之后启航，大约三十天到达奥克利斯（Ocelis）。从那里顺风，四十天就能到达印度最近的商埠穆齐里斯（Muziris）\cite{pliny}。回程要在埃及的提比月（Tybis）之初、也就是罗马历的十二月启航，最晚不过埃及梅基尔月（Mechir）的第六天。那一天相当于罗马历一月的伊都斯日（ides），即一月中旬。普林尼说，这样就能在一年之内往返\cite{pliny}。夏天去，冬天回，一年一个来回，日程本身就说明，这里的风随季节转向。路上并不太平：“现在每年都有船去印度，船上载着成队的弓箭手，因为那一带海上海盗猖獗。”\cite{pliny}

这股风有一个名字。普林尼说，从阿拉伯半岛的叙阿格鲁斯角（Syagrus）到印度的帕塔莱（Patale），航程一千三百三十五英里，最便利的是借助一种西风。那里管这种风叫希帕洛斯（Hippalus）\cite{pliny}。《厄立特里亚海周航记》（Periplus Maris Erythraei，下称《周航记》）却说希帕洛斯是一个人。他是“第一个通过观察港口的方位与海的形状，发现了横渡大洋的直航航线的领航员”。《周航记》接着说，那一带的西南季风似乎也叫希帕洛斯；英译者在方括号里补了一句，说风是因这位领航员得名的\cite{periplus}。普林尼只说这是风的名字，照《周航记》英译的读法，则是先有人，后有风的名字。《周航记》也说，从埃及去印度这一带，合适的季节大约在七月，即厄庇菲月（Epiphi）\cite{periplus}，与普林尼说的仲夏一致。

两部书的年代都不确定。《周航记》不知作者。Darley（2014）认为，把文本同纳巴泰王表、纳哈帕纳（Nahapana）可能的在位年代对起来，最好的重合当然落在公元 40—70 年这几十年。把它定在 3 世纪的说法，没有得到多少支持\cite{darley2014}。普林尼的《自然史》则完成于公元 79 年以前\cite{darley2014}。下文说 1 世纪，指的就是这两部书的时代。

\subsubsection{港口与货物}
《周航记》讲了印度西海岸的两组港口。北边的巴利加萨（Barygaza，英译本写作 Barugaza）进口酒，主要是意大利的。它还进口金币与银币，换成当地的钱有利可图；希腊原文这里写的就是“第纳尔”（δηνάριον）\cite{periplus}。当地还用着旧钱。直到《周航记》的时代，刻着阿波罗多托斯（Apollodotos）与米南德（Menander）希腊铭文的旧德拉克马，还在巴利加萨流通\cite{periplus}。这两位是早已不在的印度—希腊国王，他们的钱还在用。

南边是穆齐里斯与内尔肯达（Nelkynda）。来这些港口的船都很大，因为船上装的胡椒与蒌叶又多又占地方。运进来的首先是“大量的钱币”，然后是少量的素布，以及白玻璃、铜或黄铜、锡、铅。酒不多，大约与巴利加萨一样多；船上带的粮食只供船员吃，因为商人不卖它\cite{periplus}。运出去的头一项是大量的胡椒，只产在这些商埠中的一处，叫作科托纳拉（Kottonara）胡椒\cite{periplus}。普林尼说，那个产胡椒的地区用整棵树挖成的小船，把胡椒运到巴拉刻（Barace）\cite{pliny}。

至少在南边的港口，钱币排在进口货的头一项，布与酒都不多。普林尼的抱怨，就从这里来。这笔账该怎样记，是\secref{sec:india-bop}的问题。

\subsubsection{汉朝的南海航路}
《汉书·地理志》记下了另一条路。从日南、徐闻、合浦出海，“船行可五月，有都元国，又船行可四月，有邑卢没国；又船行可二十余日，有谌离国；步行可十余日，有夫甘都卢国。自夫甘都卢国船行可二月余，有黄支国”（图~\ref{fig:india_route}(b)）。去的是官府的人。黄门属下的译长与应募的人一起出海，带着黄金与各色丝织品，去买明珠、璧流离、奇石异物。路上靠别人的船：“所至国皆禀食为耦，蛮夷贾船，转送致之。”\cite{hanshu}

风险很大：“亦利交易，剽杀人。又苦逢风波溺死，不者数年来还。”西汉末年也派过人去。平帝元始年间（元始元年为 1 年），王莽辅政，想炫耀威德，厚赠黄支王，让他派使者来献活犀牛。回程“自黄支船行可八月，到皮宗；船行可二月，到日南、象林界云”，维基文库本把后一段的二月写作八月。黄支以南还有已程不国，汉朝的译使到那里就回来了\cite{hanshu}。黄支在哪里，这些史料没有说。

\begin{example}{两条路有多远}\label{ex:india-route}
（1）按普林尼的日程，从亚历山大里亚到穆齐里斯要多少天？尼罗河上与沙漠里，每天各走多少英里？（2）按《汉书》，从日南等地出发到黄支，至少要多少天（一个月按 30 天，“余”取下限）？回程呢？（3）汉朝的去程是罗马的几倍？为什么不能直接比？
\solution （1）$12+12+30+40=94$ 天，三个多月；从贝勒尼基算起是 70 天。尼罗河上每天 $308/12=25.7$ 英里，沙漠里每天 $257/12=21.4$ 英里。（2）去程 $5+4+2=11$ 个月，加 20 天船行与 10 天步行，至少 $330+30=360$ 天，约一年；回程 $8+2=10$ 个月，300 天，按维基文库本是 $8+8=16$ 个月，480 天。（3）汉朝的去程是罗马的 $360/94=3.8$ 倍。汉朝的使者沿着海岸一站一站地走，中间还要步行，靠的是“蛮夷贾船，转送致之”；罗马的船借季风横渡大洋，从奥克利斯到穆齐里斯只要 40 天。两条路也不能直接比：起点、终点与路程都不同，黄支在哪里史料又没有说，3.8 倍只是两部书的日程之比；何况“不者数年来还”，汉朝的使者往往几年才回得来。
\end{example}

\fig{india_route}{两条去印度的路。(a) 看罗马的去程怎样分成四段：按普林尼的日程，尼罗河 12 天、沙漠 12 天、红海与阿拉伯沿岸 30 天、横渡大洋 40 天，共 94 天；回程在十二月到一月中旬启航，一年之内往返。(b) 看汉朝的路有多长：按《汉书·地理志》，去程至少 360 天（一个月按 30 天，“余”取下限），回程 300 天；虚线是维基文库本把回程第二段写作“八月”时的 480 天。数据：普林尼《自然史》第六卷（Bostock 与 Riley 英译），《汉书·地理志下》（殆知阁数字文本，另与维基文库本比对）。}{fig:india_route}

\section{打印币与《政事论》}\label{sec:india-artha}

罗马的钱运到印度，遇到的是另一种钱，也是另一种国家。古代印度的钱是什么样的？《政事论》里的国家怎样管钱、盐、矿与商业？

\subsubsection{打印币}
在印度各地出土的最早的钱，是一种打着几个印戳的方形银片。Cunningham（1891）说，它们一面或两面打着几个印戳，因此得了一个描述性的名字，叫打印币（punch-marked）\cite{cunningham1891}。做法一看钱就明白。先把一块银子打成一张薄片，大约一枚英国先令硬币那么厚，再切成窄条；每一条切成重量相同的小块，约 56 格令，较重的坯子再切下一两个角来调整重量\cite{cunningham1891}。格令（grain）是英制的重量单位，1 格令合 0.06479891 克。56 格令约合 3.63 克。

这种钱用得很久。菩提伽耶（Bodh Gaya）佛座下的一批供物里有三枚打印币，共重 111 格令，平均只有 37 格令。Cunningham 说，这批供物是公元 150 年前后、印度—斯基泰国王胡毗色伽（Huvishka）在位时放进去的，可见打印币那时仍在流通。从 56 格令磨到 37 格令，少了 34\%。全印度八百多枚这种钱，平均重量在 47 格令以上。Cunningham 愿意接受这三枚在大约六百年的流通中损失了 19 格令，但把它当作一个非常特殊的例子\cite{cunningham1891}。这三枚每世纪损失 3 格令以上，一般的打印币每世纪不过一格令半。

打印币从什么时候开始铸，年代没有定论。Cunningham 说，它们“在佛陀的时代，即公元前 6 世纪，肯定已经通行；但我看不出它们为什么不能早到公元前 1000 年”\cite{cunningham1891}。这是一位 19 世纪学者的看法。

\subsubsection{一部年代不定的手册}
《政事论》（Arthaśāstra）是一部讲治国的书，本章据 Shamasastry 1915 年的英译本。它自称出自推翻难陀王朝的人之手：“这部论著，是那个不能容忍（暴政）、迅速把经典、兵法与落入难陀王之手的大地拯救出来的人写成的。”\cite{arthashastra}

它的年代有两种看法。Fleet 在 1914 年为译本写的序言里说，这部书自称写于公元前 321—前 296 年，古朴的文体也与这一说法相合。他还说，有见识的学者似乎都同意，现存的文本也许不是考底利耶（Kauṭilya）一字不差写下的，但基本上就是他在那个时期写的书。人们从希腊作家麦加斯梯尼（Megasthenes）的一些话里知道的很多东西，这部书不仅能印证，还能加以扩充\cite{fleet1914}。Fleet 也交代了文本的来源，译者只有一部抄本和一部不完整的注释\cite{fleet1914}。

今天的看法不同。Francis 在评论 Olivelle（2013）的新译本时，转述了 Olivelle 的结论。考底利耶的那一版，大概写成于公元 50 年到 125 年之间。175 年到 300 年之间的一次改编，使这部书有了论（śāstra）的形式。把考底利耶等同于旃陀罗笈多·孔雀（Candragupta Maurya）的大臣阇那迦（Cāṇakya），“是笈多时代的一种虚构”。Francis 认为这些论证大体可信，但哪些段落属于哪一层，讨论并没有因此结束\cite{francis2013}。

本章把《政事论》当作一部年代不确定的规范性文本。它说的是国家应当怎样做，不是哪一个王朝实际做了什么。按 Olivelle 的年代，它与普林尼和《周航记》相去不远；按 Fleet，它要早三四百年。下面引的条文，都按这个限度来读。

\subsubsection{铸币与验币}
《政事论》的钱叫帕纳（paṇa）。书里规定：“铸币长官（lakshnádhyakshah）要铸造银币（rúpyarúpa），用四份铜，加十六分之一份（másha）的 tikshna、trapu、sisa 或 anjana 中的任一种金属。钱有一帕纳、半帕纳、四分之一帕纳与八分之一帕纳。”\cite{arthashastra}英译没有说铜与那种金属合在多少银里。同书讲度量衡的一章说，16 银马沙（másha）或 20 粒 saibya 种子合 1 达拉纳（dharana）\cite{arthashastra}。马沙与达拉纳都是重量单位。如果一枚钱按 16 份计，四份铜、一份别的金属，银就是 11 份，含银 68.75\%。这只是一种读法。照这个读法，它比尼禄以后的罗马第纳尔（80\%，表~\ref{tab:greek-silver}）还要低。

钱铸出来，还要验。验币官（rúpadarsaka）既管作为交换媒介的钱，也管可以缴入国库的法偿货币。向国库交钱时，要另收“8\%，叫 rúpika；5\%，叫 vyáji；每一百帕纳八分之一帕纳，叫 páríkshika（检验费）”\cite{arthashastra}。三项合起来是 13.125\%，也就是说，用钱向国家交税，要多付一成多。国家不只从铸钱上得利，也从收钱上得利。

\begin{example}{一枚帕纳}\label{ex:india-mint}
（1）按上面的读法，一枚帕纳含银多少？与 64 年以后的罗马第纳尔相比如何？（2）一个人要向国库交 100 帕纳，连同三项附加要付多少？（3）三项附加里，检验费只占多少？其余两项更像什么？
\solution （1）$11/16=68.75\%$，比第纳尔的 80\% 低 11.25 个百分点。（2）$100\times(1+0.08+0.05+0.00125)=113.125$ 帕纳。（3）检验费是每百帕纳八分之一帕纳，只占 0.125\%，是验钱的成本；8\% 与 5\% 远远超过验钱的成本，更像是一种税：国家规定了什么钱可以交进国库，也就能对交进来的钱收费。这与\chref{ch:medieval}的铸币税（式~\eqref{eq:medieval-seigniorage}）不同：铸币税是铸钱时得的，这里是收钱时得的。
\end{example}

\subsubsection{盐与矿}
《政事论》要国家握住矿与盐。采矿与矿产品的买卖，都归国家专营。盐的长官要按时收取钱租（prakraya）和应归国家的那份盐。卖掉这些盐，不仅要收回它的价值（múlyam），还要加收 5\% 的溢价（vyájím），都收现钱（rúpa）。进口的盐（ágantulavanam）要向国王交六分之一。在盐里掺假的，处以最重的罚金；除隐士（vánaprastha）以外，未经许可制盐的也一样\cite{arthashastra}。这与\chref{ch:qinhan}的盐铁专卖（\secref{sec:qinhan-monopoly}）很像，都是国家握住必需品的供给，把税加在价格里。

\subsubsection{商业长官}
商业长官要弄清各种商品有没有需求、价格是涨是跌，不论它们是陆上还是水里的产品，从陆路还是水路运来。他要把分散在各处的商品集中起来，“抬高价格。抬高的价格被接受以后，再宣布另一个价格”，但要避免那种会伤害百姓的高额利润\cite{arthashastra}。另一章说得更具体，商业长官要为本地商品在规定的价格之上定 5\% 的利润，外国商品定 10\%\cite{arthashastra}。

对外贸易，《政事论》的态度与普林尼正好相反。它要优待进口，对进口外国商品的航海者（návika）与商人减免贸易税，让他们能得到一些利润\cite{arthashastra}。可是讲通行税的一章又说：“进口商品要按价值的五分之一交通行税。”\cite{arthashastra}同一章接着给许多货物定了更低的税率。花果、蔬菜、干鱼等收六分之一，酒、象牙等收十分之一或十五分之一，布、粮食、油、糖等收二十分之一或二十五分之一\cite{arthashastra}。几条放在一起读，大概五分之一是一般的税率，许多货物另有低一些的税率，要鼓励的进口再给予减免。英译本身没有说清它们的关系。

罗马在亚历山大里亚对东方的货物收四分之一（\secref{sec:india-cargo}），比五分之一还高。把国王的货运到外国去卖时，《政事论》要长官算一算，付完外国国王的通行税等各项费用以后，“还有没有利润的余地”\cite{arthashastra}。它担心的是有没有利润，不是金银流出流进。

\subsubsection{仓储与荒年}
国家还要有储备。这样收来的储备，一半要留着防百姓遭灾，只用另一半，旧的储备要用新的替换\cite{arthashastra}。荒年时，国王可以把自己储存的粮食或富人囤积的收入分给百姓，以示恩惠\cite{arthashastra}。这与\chref{ch:qinhan}的常平仓（\secref{sec:qinhan-granary}）同中有异。两者都要储备、都要轮换；但常平仓用一个买价与一个卖价框住粮价，《政事论》只讲留多少，没有讲按什么价钱买卖。开支也有规矩，国王供养臣仆的费用，应定在全部收入的四分之一以下\cite{arthashastra}。

\begin{compare}{《政事论》与《盐铁论》}{《政事论》}{《盐铁论》（\chref{ch:qinhan}）}
\cmprow{文本}{一部规范性的手册，讲国家应当怎样做；年代不定}{一场辩论的记录（公元前 81 年），赞成与反对两方都有}
\cmprow{专营}{采矿与矿产品的买卖；盐收租、收份额，加收 5\%}{盐、铁、酒}
\cmprow{价格}{商业长官集中商品、抬高价格，利润定在 5\%（本地）与 10\%（外国），但要避免伤害百姓}{平准：便宜时买进，贵时卖出}
\cmprow{对外贸易}{优待进口商，可减免贸易税；进口通行税一般是五分之一，进口盐六分之一}{本章的引文不涉及}
\cmprow{储备}{一半留着防灾，旧的用新的替换}{常平仓：谷贱时买，谷贵时卖}
\cmpsum{两部书都让国家握住必需品与价格：《政事论》说的是应当怎样做，《盐铁论》记下的是做了以后的争论。}
\end{compare}

\section{按风险分档的利率}\label{sec:india-interest}

在《政事论》里，同样借一百帕纳，有人每月付一又四分之一帕纳的利息，有人要付二十帕纳。利率为什么分成四档？它们说明债主怎样估计风险？

\subsubsection{四档利率}
《政事论》规定：“每月每百帕纳一又四分之一帕纳的利息是公正的。每月每百帕纳五帕纳是商业利息（vyávaháriki）。每月每百帕纳十帕纳，通行于林中。每月每百帕纳二十帕纳，通行于海上商人（sámudránám）。”\cite{arthashastra}超过上述利率、或使别人超过的，处以初级罚金。粮食借贷另有上限，丰收时节，粮食的利息按钱折算，不得超过一半\cite{arthashastra}。

按单利乘以 12，“公正”的利率一年是 15\%，商业利率是 60\%（图~\ref{fig:india_rates}(a)）。林中是 120\%，海上是 240\%。同一部书，借给商人的利率是公正利率的 4 倍，借给海上商人的是 16 倍。

为什么差这么多？\chref{ch:africa}的式~\eqref{eq:africa-risk} 给出一种读法。如果借款人只有 $p$ 的概率能还上，还不上就什么也收不回，一个风险中性（只看预期收益）的债主要求 $p(1+i)=1+r_f$，即 $p=(1+r_f)/(1+i)$。式中 $p$ 是债主估计的还款概率，$i$ 是借款利率，$r_f$ 是几乎没有风险的收益率。把“公正”的 15\% 当作 $r_f$（教学假设），借款期限取一年。商业利率对应 $p=1.15/1.60=0.72$，林中 $1.15/2.20=0.52$，海上 $1.15/3.40=0.34$（图~\ref{fig:india_rates}(b)）。按这个读法，借钱给海上商人，债主预期只有三分之一的借款能连本带利收回。

这个读法有三个限度。第一，这些是法定的利率，不是成交的利率。文本说超过要罚，没有说实际收多少；同一节里“林中”与“海上”的说法，也像是在承认当地通行的价钱，而不是在规定它。第二，式~\eqref{eq:africa-risk} 假定违约时什么也收不回；有抵押时，$p$ 应读作债主预期能收回的比例（\chref{ch:africa}）。第三，公正利率本身也含风险，$r_f$ 取得越低，算出的 $p$ 越低。

\subsubsection{海事贷款}
罗马这一边也有这样的借款，叫作\textbf{海事贷款}，即商人为一趟航行借钱，拿船或货作担保。Casson（1990）认为，纸草背后的那笔借款必定是一笔海事贷款，借款人有了它才买得到这些货。合同在哪里订的，看法不同，见下面的案例框。这一笔的担保不是船，而是要交 25\% 进口关税的货物。几乎可以肯定，借款人是拿借来的钱买货，再把货押给债主，正像几个世纪以前德摩斯梯尼时代的商人那样\cite{casson1990}。关键在于，“如果船沉了，担保就毫无意义”\cite{casson1990}。船沉了，债主什么也收不回，所以海事贷款的利率里，包含了船沉的概率。\chref{ch:islamic}的康曼达（\secref{sec:islamic-partnership}）是另一种分担风险的办法，它是一方出钱、另一方带着钱货出海经营、赚了按约定分利的合伙。海事贷款用高利息补偿债主，康曼达让出钱的人分享利润，也分担亏损。

\begin{example}{海上利率里有什么}\label{ex:india-risk}
（1）把《政事论》的四档利率折成年率（单利）。（2）取 $r_f=15\%$，按式~\eqref{eq:africa-risk} 求每一档的 $p$。（3）假定一笔一年期的海上贷款以货物作抵押：船平安到达，债主一定能收回本息；船沉了，什么也收不回。如果 240\% 的利率只补偿船沉的风险，船沉的概率是多少？这可信吗？
\solution （1）$1.25\times12=15\%$，$5\times12=60\%$，$10\times12=120\%$，$20\times12=240\%$。（2）$1.15/1.15=1$，$1.15/1.60=0.72$，$1.15/2.20=0.52$，$1.15/3.40=0.34$。（3）船平安到达的概率就是 $p=0.34$，船沉的概率是 $0.66$，每三条船沉两条。普林尼说每年都有船去印度，斯特拉波说有一百二十条，这样的沉船率不可信。所以这一档利率里还有别的东西：借款人本身还不上的风险、比一年更长或更短的期限，或者它只是一个法定的上限，不是成交的价钱。式~\eqref{eq:africa-risk} 能把利率换成概率，却不能告诉我们利率里各是什么。
\end{example}

\fig{india_rates}{《政事论》的四档利率。(a) 看利率怎样随借款人变：每月每百帕纳 1.25、5、10、20 帕纳，单利折成年率是 15\%、60\%、120\% 与 240\%。(b) 看这些利率意味着什么：把“公正”的 15\% 当作几乎没有风险的收益率（教学假设），按式~\eqref{eq:africa-risk}，债主估计的还款概率从 1 降到 0.34。数据：《政事论》第三卷第 11 章（Shamasastry 英译，维基文库）。}{fig:india_rates}

\section{一张纸草上的一船货}\label{sec:india-cargo}

普林尼的一亿是整个帝国一年的数，纸草上记的却是一条船的账。一条从印度回来的船，装的货值多少？国家从中拿走多少？

\subsubsection{纸草的两面}
Casson（1990）解读的这张纸草，两面大致同时，都写于 2 世纪中期\cite{casson1990}。正面是协议的一部分，讲货物在红海边卸下以后怎样运到亚历山大里亚。借款人说：“我会在看守与保护之下，把（货物）经沙漠运到科普托斯收取税款的公共仓库”，然后顺流运到亚历山大里亚收四分之一关税的仓库\cite{casson1990}。背面算的是货物的价值，用来交这四分之一的关税。

货单上的塔兰特有两种用法。后面跟着弥那（mina）的是重量，弥那在这里也是重量单位，每塔兰特合 60 弥那。后面跟着德拉克马的是钱，与开头一样。1 德拉克马合 6 奥波尔（obol）。这些比例与\chref{ch:greekroman}讲的希腊单位相同，用纸草自己的算术就能解出（例~\ref{ex:india-cargo}）。表~\ref{tab:india-cargo} 是背面保存下来的三批货。

\begin{table}[htbp]\centering\small
\caption{赫尔玛波隆号的货（纸草背面第二栏）}\label{tab:india-cargo}
\begin{tabularx}{\textwidth}{@{}lYrr@{}}
\toprule
货物 & 计价 & 价值 & 德拉克马 \\
\midrule
恒河甘松，60 箱 & 每箱 4500 银德拉克马 & 45 塔兰特 & 270,000 \\
象牙 & 计税重量 76 塔兰特 45 弥那，每弥那 100 德拉克马；另加红海卸货港的税官（Arabarchs）取走的 $11\tfrac34$ 弥那 & 76 塔兰特 5675 德拉克马 & 461,675 \\
布，54 匹 & 计税重量 12 塔兰特 27 弥那，每弥那 70 德拉克马；另加税官取走的 $22\tfrac34$ 弥那 & 8 塔兰特 5882 德拉克马 3 奥波尔 & 53,882.5 \\
\midrule
保存下来的三批 & & & 785,557.5 \\
六批合计 & “以银计” & 1154 塔兰特 2852 德拉克马 & 6,926,852 \\
\bottomrule
\end{tabularx}
\datasource{Casson（1990），第 200—202 页（据 Thür 对背面第二栏的解读译出）。塔兰特后面跟着弥那的是重量，跟着德拉克马的是钱；计税重量是亚历山大里亚的税官按惯例给商人减让一点以后的重量；1 塔兰特（重量）= 60 弥那，1 塔兰特（钱）= 6000 德拉克马，1 德拉克马 = 6 奥波尔，都由纸草自己的算术解出（例~\ref{ex:india-cargo}）。最后一栏由本章折算，见 \texttt{checks/check\_ch09.py}。}
\end{table}

\begin{case}{2 世纪中期 · 穆齐里斯—亚历山大里亚 · 赫尔玛波隆号}
借款人借钱买货，合同在哪里订的，两位研究者看法不同。据 Casson 的转述，Thür 认为借款合同是出航以前在亚历山大里亚订的，借款与担保各写一份文书；纸草上提到穆齐里斯的那个短语，Thür 理解为去穆齐里斯的航行。Casson 按这个短语通常的意思来读，倾向于认为，“原来的合同是在穆齐里斯订的”，“借贷双方或其中一方，很可能是住在那里的外国侨民”\cite{casson1990}。货物运到埃及以后，如果到期不还，债主可以处置这批货。交了四分之一的关税，剩下的四分之三，他可以运到想去的地方，卖掉、再抵押、转让给别人，随他的意\cite{casson1990}。

背面只保存了三批货，另外三批一定写在缺失的那部分里\cite{casson1990}。三批合计 78.6 万德拉克马，只占全船 692.7 万的 11.3\%。缺失的三批值 614.1 万。Casson 在注里说，这笔钱差不多等于修建亚历山大·特罗阿斯（Alexandria Troas）引水道的花费，那是 700 万德拉克马，其中赫罗狄斯·阿提库斯出 400 万，哈德良出 300 万。一条普通大小的罗马商船，载重 340 吨，能装 11000 多塔兰特这样的货。Casson 说，“它们是名副其实的宝船”\cite{casson1990}。

国家得到多少？按四分之一，这条船交的关税是 173.2 万德拉克马。斯特拉波说，亚历山大里亚收双份的关税，一份来自进口，一份来自出口；最贵重的货物课税最重，因为亚历山大里亚实际上垄断了这门贸易\cite{strabo}。东方的奢侈品，是罗马国库的一个税源。
\end{case}

\begin{example}{纸草上的算术}\label{ex:india-cargo}
（1）设一塔兰特（重量）合 $W$ 弥那，一塔兰特（钱）合 $T$ 德拉克马。用象牙与布两批货的算术解出 $W$ 与 $T$；再用布的余数求一德拉克马合多少奥波尔。（2）保存下来的三批货共值多少？占全船的多少？（3）全船要交多少关税？合多少塞斯特斯？（4）全船的价值是普林尼那一亿的多少？是印度单独的五千万（读法见\secref{sec:india-bop}）的多少？
\solution （1）象牙按 76 塔兰特 45 弥那、每弥那 100 德拉克马计价，纸草写作 76 塔兰特 4500 德拉克马：$(76W+45)\times100=76T+4500$，所以 $T=100W$。布按 12 塔兰特 27 弥那、每弥那 70 德拉克马计价，纸草写作 8 塔兰特 4290 德拉克马：$(12W+27)\times70=8T+4290$；代入 $T=100W$，得 $840W+1890=800W+4290$，$W=60$，$T=6000$。所以象牙的计税重量是 $76\times60+45=4605$ 弥那，值 460,500 德拉克马，与\chref{ch:greekroman}的换算（1 塔兰特 = 60 弥那 = 6000 德拉克马）一致。布的余数 $22\tfrac34\times70=1592.5$ 德拉克马，纸草写作 1592 德拉克马 3 奥波尔，所以半个德拉克马是 3 奥波尔，1 德拉克马 $=6$ 奥波尔。（2）$270000+461675+53882.5=785557.5$ 德拉克马，占 $785557.5/6926852=11.3\%$；全船 $1154\times6000+2852=6926852$ 德拉克马。（3）$6926852/4=1731713$ 德拉克马，按 1 德拉克马合 1 塞斯特斯，是 173.2 万塞斯特斯。（4）$692.7/10000=6.9\%$；$692.7/5000=13.9\%$。
\end{example}

\section{古代的“国际收支”}\label{sec:india-bop}

罗马的金银流向东方，按今天的账户来记，是不是一个经常账户逆差？要回答，先得弄清普林尼到底说了多少，再看这笔钱在罗马的经济里有多大。

\subsubsection{普林尼的两个数}
普林尼的数字有两个。一个就是开头的那一亿，按最低的估计，印度、赛里斯与阿拉伯半岛每年从帝国拿走一亿塞斯特斯\cite{pliny}。另一个只说印度。英译本是这样的：“没有哪一年印度从我们的帝国吸走的钱少于五亿五千万塞斯特斯，它运回来的货物，在我们这里卖到原价的整整一百倍。”\cite{pliny}

拉丁原文没有写出这个数字，写的是一个数码。LacusCurtius 网站上的 Teubner 本作“nullo anno minus HS·| D | imperii nostri exhauriente India et merces remittente, quae apud nos centiplicato veneant”\cite{pliny}。数码是 D，即 500，外面加了框。这个框是什么意思，同一页上有旁证。

先看横线。p. 是步（passus），一千步为一罗马里；数码上加一道横线，表示千。没有横线的“MM p.”，英译作两英里；有横线的“CCCVIIII p.”是 309 千步，英译作三百零八英里。

再看框。框表示再乘以一百，即十万。6.100 的一段里程写作“| XIII |· XXXII p.”。这是 13 个十万步加 32 个千步，合 1332 罗马里，英译作一千三百三十五英里。可见英译者读里程时，是把框读作十万的。普林尼自己说，古人没有比十万更大的数，所以到他那时，罗马人还是按十万的倍数来计。12.84 的那一亿，拉丁原文就用文字写出了这样一个倍数，“miliens centena milia sestertium”，即一千个十万塞斯特斯，英译作一亿\cite{pliny}。

按这种计法，框起来的 D 是五百个十万，即五千万塞斯特斯。Mayhoff 编的 Teubner 本，校勘记还记下了 DL 这一读法\cite{pliny}，即五千五百万。英译的五亿五千万，是 DL 的十倍。12.84 本身也是一个检验。三个地区加起来，按最低的估计是一亿；如果印度一处每年就不少于五亿五千万，三个地区的最低估计就不会只有一亿。下文读作五千万，另记五千五百万，这只是一种读法。

\subsubsection{放进国际收支表}
把罗马帝国当作“本国”，把印度、阿拉伯与更远的东方当作“外国”。记号沿用\chref{ch:external}的式~\eqref{eq:external-bop}。$\mathit{CA}$、$\mathit{KA}$、$\mathit{FA}$、$\mathit{EO}$ 依次是经常账户、资本账户、金融账户、净误差与遗漏，$\Delta A$ 与 $\Delta L$ 是对外资产与对外负债的净增加。罗马向东方出口货物 $X$（意大利的酒、玻璃、铜、锡、铅、布），进口货物 $M$（胡椒、象牙、甘松、布），差额用钱币支付，运出的钱币记作 $C$。假定罗马人不向东方借钱，东方人也不持有对罗马的债权（$\Delta L=0$），资本账户与误差都不计，那么 $X+C=M$。这笔钱币有两种记法。
\begin{keyeq}
\underbrace{\mathit{CA}=X+C-M=0}_{\text{钱币当作货物}}
\qquad\text{或}\qquad
\underbrace{\mathit{CA}=X-M=-C=\Delta A=\mathit{FA}}_{\text{钱币当作罗马的金融资产}} .\label{eq:india-coin}
\end{keyeq}
两种记法里，$X$ 与 $M$ 都一样，不同的只是 $C$ 记在哪里。按第一种记法，钱币与酒一样是出口的货物，经常账户是平衡的，没有任何借贷。按第二种记法，钱币是罗马的财富，运出去就是对外资产的减少；经常账户有一个大小为 $C$ 的逆差，用减少的资产来弥补。普林尼说印度在“吸干”（exhauriente）帝国，他的看法属于第二种。

今天的规则在两者之间划了一条线。按美国经济分析局（BEA）的手册，货币黄金是货币当局拥有所有权、作为储备资产持有的黄金；非货币的实物黄金不是金融资产，归为货物\cite{bea2026iea}。罗马商人手里的金币，不是哪个货币当局的储备，按今天的规则应当记作货物，即第一种记法。手册还写明，按国际收支的口径，非货币黄金包括“非纪念性的金币”\cite{bea2026iea}。手册这里讲的是黄金，本章把银币也按同样的道理处理。

还有第三种可能，要放开 $\Delta L=0$ 的假定。手册把通货定义为由中央银行或政府发行或授权、面值固定的纸币与硬币，列在债务工具里\cite{bea2026iea}。今天的硬币，是发行者的负债。第纳尔是罗马政府发行的，在罗马境内有法定的面值。如果印度人按这个面值收下它，再拿它向罗马人买东西，它就是印度对罗马的债权。这时经常账户仍有大小为 $C$ 的逆差，只是用增加的负债（$\Delta L=C$）来弥补。可是在印度没有人保证它的面值，这种读法不合适。

印度人是把罗马钱当作货物看的。《周航记》说金币与银币换成当地的钱有利可图，换钱的人看的是这些钱币本身值多少当地的钱\cite{periplus}。Darley 讲晚期罗马的钱时说，把印度的罗马钱当作贸易不平衡的证据，“没有看到要点：这些钱币在印度市场上是一种有价值的商品，凭它们的外观，作价高于所含的金属”\cite{darley2014}。

这并不是说外流无关紧要。钱币运出去，罗马国内的钱就少了一点。钱少了，物价可能低一点，这是\chref{ch:bullion}要讲的价格—铸币流动机制。问题是这一点有多大。

\subsubsection{有多大}
罗马的 GDP 没有统计，只有估计。Scheidel（2019）估计，2 世纪罗马帝国的 GDP 略少于 200 亿塞斯特斯。他在注里说，他与 Friesen 合作计算时用的是 180 亿，这是一个更宽的可能区间的中点。他也说，GDP 不可能低很多，否则社会就会普遍贫困，而考古记录没有反映这一点\cite{scheidel2019}。人口，Scheidel 为简单起见取 7000 万\cite{scheidel2019}，人均约 257 塞斯特斯。

放到罗马的经济里，普林尼的外流是一个很小的数。以 180 亿为分母，印度的五千万是 GDP 的 0.28\%，五千五百万是 0.31\%（图~\ref{fig:india_scale}）。三个地区的一亿是 0.56\%。即使是英译的五亿五千万，也只有 3.1\%，还不到美国 2025 年经常账户逆差占 GDP 的 3.8\%（\chref{ch:external}）\cite{wdi2026}。赫尔玛波隆号一条船的货，是 GDP 的 0.04\%。用 2 世纪的 GDP 去量 1 世纪的数字，只能看数量级，但数量级是清楚的。

另一把尺子是钱的存量。Duncan-Jones 估计，到 2 世纪中期，罗马帝国全部钱币的价值约为 200 亿塞斯特斯。Scheidel 转引了这个数，但认为它很可能偏高；他说，一个价值相当于全年 GDP 的铸币存量，会让罗马成为一个非常罕见的例子\cite{scheidel2019}。按 200 亿计，每年五千万是存量的 0.25\%。如果一百年里年年如此，就是 50 亿，相当于存量的四分之一，这还没有算新开采的金银。存量越小，这个比例越大。外流对罗马的影响如果有，是在钱的存量上慢慢积累的，不是一年之内的“逆差”。

\begin{example}{一百倍是多少}\label{ex:india-scale}
普林尼说，印度的货在罗马“卖到原价的整整一百倍”。（1）如果罗马每年付给印度五千万塞斯特斯，按一百倍，这些货在罗马值多少？占 GDP（180 亿）的多少？（2）这可信吗？（3）按人均 GDP，每个罗马人每年要为印度货花多少？
\solution （1）$5000\text{ 万}\times100=50$ 亿塞斯特斯，是 GDP 的 $50/180=27.8\%$。（2）不可信：一个人均 257 塞斯特斯、普遍不富裕的社会，不可能把四分之一以上的收入花在胡椒、象牙与宝石上。“一百倍”更像是修辞，与“按最低的估计”一样，是在强调这笔钱之大。（3）$50\text{ 亿}/7000\text{ 万}=71$ 塞斯特斯，是人均 GDP 的 27.8\%。
\end{example}

\fig{india_scale}{普林尼的外流有多大。看每一根柱子占罗马 GDP 的比例（对数横轴；GDP 按 Scheidel 与 Friesen 计算用的 180 亿塞斯特斯）：赫尔玛波隆号一条船的关税与货值；印度的五千万（本章的读法）与五千五百万（校勘记里的读法）；三个地区的一亿；英译的五亿五千万；以及按“一百倍”算出的印度货在罗马的价值；灰色的两根是本章认为不可信的数字（\secref{sec:india-test}）。虚线是美国 2025 年的经常账户逆差，GDP 的 3.8\%。数据：普林尼《自然史》第六、十二卷，Casson（1990），Scheidel（2019），世界银行 WDI（\chref{ch:external}）。}{fig:india_scale}

\begin{compare}{贵金属外流与经常账户逆差}{罗马的金银外流}{今天的经常账户逆差（\chref{ch:external}）}
\cmprow{看的是什么}{钱币（金属）流出帝国}{货物、服务与收入的收支差额}
\cmprow{怎样弥补}{用钱币支付，没有借贷}{向外国借钱，或减少对外资产}
\cmprow{钱币算什么}{在印度是一种商品；按今天的规则，与非货币黄金一样记作货物}{储备黄金是金融资产；政府发行、面值固定的硬币是通货，算发行者的负债}
\cmprow{有多大}{普林尼的五千万到一亿，GDP 的 0.3\%—0.6\%}{美国 2025 年，GDP 的 3.8\%}
\cmprow{担心什么}{帝国的财富被“吸干”，钱花在奢侈与女人身上}{外债能不能持续，汇率会不会大跌}
\cmpsum{普林尼看到的是金属的流出，今天的账户看的是借贷；用钱币付账而不借钱，按今天的规则，经常账户可以是平衡的。}
\end{compare}

\begin{pitfall}[金银流出，国家就变穷]
罗马把金银运到印度，换回胡椒与象牙。付出去的金银是买东西的钱，不是丢掉的钱。罗马人得到了胡椒，印度人得到了金银，双方都认为换得值。金银流出使罗马的钱少了一点，可能使物价低一点（\chref{ch:bullion}），却不使罗马人更穷，除非胡椒不值那个价钱，而那是一个关于奢侈的道德问题，不是一个关于账户的问题。把金银的流出当作国家的损失，一千多年后东印度公司运银去东方时还有人这样抗议。替公司辩护的孟说，运出去的白银只要用来做贸易，买回的货再卖出去，就会带回更多的金银，所以不是损失；但孟仍然用流入的金银来量国家的财富（\chref{ch:mercantilism}的\secref{sec:merc-mun}）。《政事论》的商业长官问的是“还有没有利润的余地”。
\end{pitfall}

\section{罗马钱在印度}\label{sec:india-coins}

账怎样记是一回事，印度的地下埋着什么，是另一回事。在印度出土的罗马钱，是什么时候、什么金属的？它们说明了什么？

\subsubsection{按世纪数}
Darley（2014）把发表过的、在印度出土的罗马与早期拜占庭钱币，按发行者的世纪汇总（图~\ref{fig:india_coins}(a)）。1 世纪以前，银币（AR）7 枚，金币（AV）6 枚。1 世纪最多，银币 4174 枚，金币 756 枚。2 世纪有金币 378 枚。3 世纪只有金币 12 枚。4 世纪有金币 1 枚、铜币（AE）616 枚。5 世纪有金币 84 枚、铜币 101 枚。6 世纪有金币 21 枚。7 世纪只有金币 2 枚。Darley 提醒，这些数字只是提示性的\cite{darley2014}。合计 6158 枚，1 世纪的占 80.1\%。银币共 4181 枚，全是 1 世纪以前与 1 世纪的。Darley 的图 5.7 给的合计又不一样，罗马钱 6446 枚，其中银币 4792 枚\cite{darley2014}。这些数只能看大势。

这些数按发行者计，不是按钱到印度的时间计。一批晚埋的窖藏，可能全是旧钱。第纳尔的大窖藏可以有上千枚。Darley 转引了 Turner（1989）列出的大窖藏。阿肯帕莱（Akenpalle）约 1531 枚，布迪纳塔姆（Budinatham）1398 枚，韦拉卢尔（Vellalur）522 枚\cite{darley2014}。

\subsubsection{旧的银币，后来的金币}
1 世纪的银币是谁发行的？按 Darley 的附录 6 加总，能认出发行者的银币有 2497 枚（图~\ref{fig:india_coins}(b)）。其中奥古斯都与提比略的有 2479 枚，占 99.3\%，包括分不清是两人中哪一位的。尼禄的只有 4 枚。同在尼禄之前的卡利古拉（Caligula）与克劳狄乌斯（Claudius），两人的银币合起来却只有 14 枚。罗马的金币叫奥里斯（aureus），据普林尼英译者的注，1 枚合 25 第纳尔\cite{pliny}。金币在印度一直延续下去。克劳狄乌斯的金币有 215 枚，尼禄的有 193 枚。2 世纪还有哈德良的 100 枚、安敦宁·毕尤的 182 枚。

\chref{ch:greekroman}讲过，第纳尔的含银在 64 年尼禄改革以前约 90\%，改革后降到 80\%（表~\ref{tab:greek-silver}）。印度的窖藏几乎只要改革以前的旧银币。Darley 认为，说印度洋贸易不要白银是站不住的，因为至少到 2 世纪晚期，罗马与印度的贸易还在用第纳尔。他接着说，印度出土的第纳尔减少，“与银币含银的不断降低同时发生，这强烈地提示：罗马对印度的贸易改用金币，是因为对罗马银币成色的不信任，而不是因为对白银本身失去了兴趣”\cite{darley2014}。两句合起来读，贸易用的还是第纳尔，出土的却几乎全是尼禄以前的，所以运去的多半是旧钱。《周航记》说的旧德拉克马也是这样，印度—希腊国王早已不在，他们的钱还在巴利加萨流通\cite{periplus}。

用\chref{ch:greekroman}的式~\eqref{eq:greek-value} 来读。式中 $c$ 是钱所含的银，$\bar v$ 是法定价值，$\tau$ 是信任。在罗马境内，皇帝的法令撑着 $\bar v$，第纳尔可以比它的银子值钱。到了印度，没有人保证 $\bar v$，最自然的读法是 $\tau=0$，钱按所含的银作价，$v=c$。可是 $c$ 要知道才行。改革以前的旧钱成色划一，看一眼就知道含多少银。64 年以后新旧两种标准并用（表~\ref{tab:greek-silver}），外观差不多，成色却不同，要知道一枚新钱的 $c$，就得重新验。

Darley 说晚期罗马的钱币在印度“凭它们的外观”作价高于所含的金属。按这个读法，外观值钱，是因为它能说明里面有多少金属。这也合乎\chref{ch:medieval}讲的劣币驱逐良币，即面值相同的两种钱，成色好的被收起或运走，留在市面上的是成色差的。在罗马境内，新旧第纳尔按同样的面值流通，含银多的旧钱留在国内并不多得什么，运到按成色作价的印度却更值钱。在印度，没有法令强迫人们按面值收钱，成色不明的新钱就不受欢迎。

\begin{intuition}[看得出成色的钱才好用]
一枚钱到了发行者的法令管不着的地方，就只是一块金属。收钱的人要知道它含多少金属：要么认得这种钱、知道它的标准，要么花工夫去验。改革以前的旧第纳尔，一看就知道是什么标准；改革以后，同样模样的钱可能是两种标准中的任何一种。对印度的商人，旧钱省了验钱的麻烦，新钱要打折扣或者重验。罗马商人于是把旧钱运去印度，把新钱留在国内按面值花掉。
\end{intuition}

这个读法有两个限度。第一，改革解释不了全部。卡利古拉与克劳狄乌斯的银币也在改革以前，出土的却只有 14 枚；各朝铸了多少银币、有多少流到东方，这些史料没有数字。第二，Darley 在评论印度人专门囤积某几种罗马钱的说法时提醒过，说某些钱币更受信任，等于假定印度商人不善于鉴定金银，而这并没有证据\cite{darley2014}。如果印度商人会验钱，新钱按成色打个折扣就是了，未必要拒收。出土的钱与成色的解释一致，却不能排除别的解释（\secref{sec:india-test}）。至于金币，Howgego 等人讨论埃及的币制时说，黄金大概经过埃及流到印度，在那里有相当数量的发现\cite{howgego2010}。

\begin{example}{出土的钱能量出流量吗}\label{ex:india-finds}
（1）按 Darley 的世纪合计，印度出土的罗马钱共多少枚？银币、金币、铜币各多少？1 世纪的占多少？（2）印度每年五千万塞斯特斯，合多少第纳尔？多少枚奥里斯？1 世纪出土的 4174 枚银币是一年流量（全用第纳尔支付）的多少？（3）出土的钱能说明什么，不能说明什么？
\solution （1）6158 枚；银币 $7+4174=4181$，金币 $6+756+378+12+1+84+21+2=1260$，铜币 $616+101=717$；1 世纪 $(4174+756)/6158=80.1\%$。（2）$5000\text{ 万}/4=1250$ 万第纳尔，或 $5000\text{ 万}/100=50$ 万枚奥里斯；$4174/1250\text{ 万}=0.033\%$。（3）出土并发表的钱，是流入的钱里极小的一部分：大部分被熔掉、改铸、遗失，或者出土了没有发表。它们能说明钱的构成（什么金属、谁发行的）与大致的时间，不能说明流量有多大，更不能用来检验普林尼的数字。
\end{example}

\fig{india_coins}{印度出土的罗马钱。(a) 看金属怎样随世纪变：按 Darley（2014）图 5.8 的注所列的合计，1 世纪以银币为主（4174 枚），以后几乎只有金币，4、5 世纪有一批铜币（卡鲁尔（Karur）另有 4120 枚铜币分不出是哪个世纪的，没有画）；对数纵轴。(b) 看 1 世纪的银币是谁的：按 Darley 附录 6 加总，能认出发行者的银币几乎全是奥古斯都与提比略的（含两人难分的 233 枚），尼禄的只有 4 枚；金币则延续下去，尼禄以后的皇帝也有。另有 1675 枚银币认不出发行者，没有画。数据：Darley（2014），附录 6 与图 5.8 的注（\texttt{data/fetch/darley2014\_coin\_finds.py}）。}{fig:india_coins}

\section{外流有多严重：三种看法}\label{sec:india-debate}

同一批史料，读出了三种看法。普林尼看到的是挥霍，Casson 看到的是宝船，Darley 则认为普林尼的数字不能拿来计算。罗马与东方的贸易，对罗马是不是一个负担？各方最有力的证据是什么？

\subsubsection{道德家：普林尼与塔西佗}
普林尼的话本身就是一种道德评判，他算的是罗马为奢侈与女人付的代价。塔西佗（Tacitus）记下了提比略在 22 年致元老院的信，抱怨的也是这个：“那种妇女特有的奢侈，为了宝石，把我们的财富转移到外族或敌国手里。”\cite{tacitus_cb}这一方最有力的理由是，帝国用金银换来的是很快就被消费掉的东西，而普林尼的数字按最低的估计就有一亿。

\subsubsection{怀疑者：Raschke 与 Darley}
Darley 认为普林尼的数字不能用来计算。他说，即使档案存在、普林尼也看到了，“他的数字是整数，又在他抨击罗马精英挥霍的长篇议论里起着修辞的作用，这应当让人在用它们做定量分析时极其谨慎”。他的论点不是说普林尼肯定没有看到政府的档案，而是说，没有证据，就不能假定他看到了\cite{darley2014}。

Darley 还引了 Raschke（1978）的话：“所有这些虚假的数字都应当抛弃。它们暗示了一种证据所不能支持的准确与精确。”他自己的论文因此避免一切绝对的价值计算\cite{darley2014}。他的结论是，伊斯兰兴起以前地中海与印度之间的贸易，“对任何一个参与的国家来说，意义都微乎其微”\cite{darley2014}。

\subsubsection{乐观者：Casson 与 Young}
另一方看重的是纸草。Casson 说，这样的船“是名副其实的宝船”\cite{casson1990}。Darley 批评 Young（2001）高估了普林尼那一亿塞斯特斯在数量上的可靠性。照 Darley 的转述，Young 承认，进入红海的印度货物要交百分之二十五的税，拿普林尼的数字去估计这笔税对罗马经济的价值是冒险的，可他仍然这样做了。Young 的理由是，如果假定普林尼用的是税收的记录，这个数的数量级大概是对的\cite{darley2014}。

这一方最有力的理由来自斯特拉波。亚历山大里亚“收双份的关税”，“最贵重的货物课税最重”\cite{strabo}。按四分之一的税率，一亿的进口值意味着每年 2500 万塞斯特斯的关税，相当于 14.4 条赫尔玛波隆号的关税。对罗马的经济，这笔贸易也许很小；对罗马的国库，它是一个看得见的税源。

\subsubsection{能确定什么}
三方的分歧，大半在于普林尼的数字能不能用。把它放在一边，能确定的有四件事。第一，船很大，货很贵，这有纸草为证。第二，罗马用钱币付了一部分账，有《周航记》与出土的钱为证。第三，国家从中收了重税，见纸草与斯特拉波。第四，即使照普林尼的数字，外流也只占 GDP 的 0.3\%—0.6\%。三种看法量的是不同的东西。普林尼担心的是帝国的财富，Darley 看到的是国家经济里的一个小数，Casson 与 Young 看到的是商人与国库手里的一大笔钱。

\section{印度、中国与罗马}\label{sec:india-china}

罗马的作者抱怨金银流走，中国的史家只能靠传闻了解大秦。从中国的史书看，这条商路上的金银往哪里流？三方对贵金属的需求有什么不同？

\subsubsection{中国史书里的大秦}
《后汉书·西域传》说大秦“以金银为钱，银钱十当金钱一。与安息、天竺交市于海中，利有十倍”。大秦王常想派使者到汉朝，安息却想用汉朝的丝绸同大秦交易，所以从中阻隔，使大秦不能自己到达\cite{houhanshu}。安息（帕提亚）夹在中间，要做丝绸的中间商。

汉朝也派人往西走过。和帝永元九年（97 年），都护班超派甘英出使大秦，到了条支。甘英到了海边要渡海，安息西界的船人对他说，“海水广大，往来者逢善风三月乃得度”；风不顺时，也有走两年的，所以出海的人都带上三年的粮食。甘英听了，就停了下来\cite{houhanshu}。

直到桓帝延熹九年（166 年），大秦王安敦的使者从日南郡边界以外来献象牙、犀角与玳瑁，两国这才第一次相通。可是“其所表贡，并无珍异，疑传者过焉”\cite{houhanshu}，史家自己也怀疑传闻夸大了。天竺西边与大秦相通，有大秦的珍物，出产各种香料、石蜜、胡椒、姜与黑盐\cite{houhanshu}。

“银钱十当金钱一”与罗马的币制对不上。按普林尼英译者的注，一枚奥里斯合 25 枚第纳尔\cite{pliny}。这可能是传闻有误，也可能说的不是罗马本国的钱。它是经过安息与天竺几道转手的消息，与普林尼的一百倍、《后汉书》的“利有十倍”一样，都说明这条商路上的人对彼此知道得很少，只知道它利润很大。

\subsubsection{三方的需求}
三方的需求不同。罗马要的是胡椒、象牙、香料与宝石，付出的是钱币与一些货物（\secref{sec:india-world}）。汉朝的使者带着黄金与丝织品出海，买的是明珠、璧流离与奇石异物\cite{hanshu}，也是奢侈品。印度收下的是金银。它有自己的银币传统，即打印币；它把罗马的钱当作商品，要成色可靠的旧银币，后来要金币。Darley 说，晚期罗马的钱币在印度凭外观作价，高于所含的金属\cite{darley2014}。

后果也不同。对罗马，外流是道德家的话题、国库的税源，在经济里却是一个小数。对汉朝，海上的路太远太险，“不者数年来还”，主要的丝绸贸易走的是经过安息的陆路。对印度，金银留了下来，埋进窖藏，或者熔掉改铸，今天出土的只是其中极小的一部分。一千多年以后，大量的白银流向中国（\chref{ch:silver}的\secref{sec:silver-china}）。那时的问题与这里一样，金属流向一个地方，是因为那里的人更看重它。

\section{数据检验与反例}\label{sec:india-test}

普林尼的书、斯特拉波的船数、一张纸草与几千枚出土的钱，出自不同的人、不同的世纪，彼此可以对账。普林尼的数字前后一致吗？纸草与斯特拉波的船数对得上普林尼吗？出土的钱能不能量出流量？

\subsubsection{普林尼对普林尼}
第一个检验在普林尼自己的书里。12.84 说三个地区每年“按最低的估计”是一亿，6.101 的英译却说印度一处每年不少于五亿五千万。两者很难都对。印度一处就有五亿五千万，三个地区的最低估计就不会只有一亿。按\secref{sec:india-bop}的读法，印度是五千万或五千五百万，约占三个地区合计的一半，两处就一致了。第二个检验是“一百倍”。五千万的货按一百倍卖，是 GDP 的 27.8\%（例~\ref{ex:india-scale}），这不可能。所以普林尼的两句话，一句的数字可以读通，另一句的倍数只能当作修辞。

\subsubsection{一条船、一百二十条船与一亿}
赫尔玛波隆号一条船的货值 692.7 万德拉克马。按斯特拉波的一百二十条船，如果这是一年的船数、每条都装这么多，一年就是 8.3 亿，是普林尼那一亿的 8.3 倍，是印度五千万的 16.6 倍。反过来，普林尼的一亿只相当于 14.4 条赫尔玛波隆号，五千万相当于 7.2 条。

对不上的原因可能有好几个。一是时代不同，斯特拉波说的是奥古斯都时代，纸草在 2 世纪中期。二是船与船不同，未必每条都装象牙与甘松。三是纸草的价值是在亚历山大里亚为交关税估的价，普林尼说的是付给印度的钱，中间隔着运费、利息、税与商人的利润。四是罗马也有货物出口，付出去的钱币只是差额。没有更多的数字，这些可能性分不开。

Casson 自己的算术也要核对。Casson 在注里说，象牙与布两批“共重 92 塔兰特，价值 528,775 德拉克马”\cite{casson1990}。按纸草上的数，两批合计是 515,557.5 德拉克马。按毛重与同样的单价算，是 528,757.5 德拉克马，与 Casson 的数差 17.5。本章用纸草上的数。

\subsubsection{出土的钱与流量}
例~\ref{ex:india-finds} 说明，1 世纪出土的 4174 枚银币，只是普林尼一年流量的 0.033\%。出土的钱不能量出流量。它们能检验的是另一件事，但能力有限。只看按发行者计的数字，分不开“对成色的不信任”与“印度不再要白银”，两种说法都预言尼禄以后的银币很少。Darley 的论证多了一条证据。贸易至少到 2 世纪晚期还在用第纳尔，出土的却几乎全是尼禄以前的，运去的多半是旧钱，这与印度不再要白银的说法不一致（\secref{sec:india-coins}）。可是同在改革以前的卡利古拉与克劳狄乌斯，银币也很少，改革不是唯一的原因。

\begin{deeper}[读这些史料时要注意什么]
\textbf{译文的译文。}普林尼、斯特拉波、《周航记》、塔西佗与《政事论》都据英译本，这些英译大多出自 19 世纪；本章的中文是英译的翻译。普林尼的外流数字另据拉丁原文，读法见\secref{sec:india-bop}，那里以原文为准，并说明这是本书的读法；英里与日数照英译，拉丁原文的里程与它略有出入（\secref{sec:india-bop}）。

\textbf{两个版本的异文。}《汉书》与《后汉书》的引文与维基文库本比对过。影响字面的不同有：黄支回程的第二段，殆知阁本作“船行可二月”，维基文库本作“船行可八月”，本章两说并存；《后汉书》的“交市”，殆知阁本作“交巿”；“瑇瑁”，殆知阁本作“檋瑁”；甘英一段讲风不顺的四个字，维基文库本作“若遇遲風”，殆知阁本的第三个字是一个错字，本章没有引这四个字。全部异文记在数据快照的说明里。

\textbf{年代。}《政事论》的年代从公元前 4 世纪到公元 3 世纪都有人主张；《周航记》大概在 40—70 年；纸草在 2 世纪中期；斯特拉波在奥古斯都时代。把它们放在一起，是因为它们讲的是同一条航线，不是因为它们同时。年号换算成公元：元始元年为 1 年，永元元年为 89 年，延熹元年为 158 年。本章的年代凡未写明公元前的，都在公元以后。

\textbf{没有考证的。}打印币从什么时候开始铸，本书没有找到可以核对的现代断代，正文只引了 Cunningham 的看法，不作结论。黄支在哪里，本书也不作考证。

\textbf{出土的钱。}Darley 的数字是发表过的钱，按发行者的世纪计。附录 6 逐条加总，与图 5.8 的注所列的合计大多一致，少数不一致（例如 2 世纪的金币，附录加总是 357 枚，注里是 378 枚）；本章的按世纪合计用注里的数，按发行者的分布用附录。论文声明未经许可不得转载，所以附录的加总不放进本书的仓库，要用时运行 \texttt{data/fetch/darley2014\_coin\_finds.py} 生成。
\end{deeper}

\section{今天还剩什么}\label{sec:india-legacy}

从普林尼的一亿到印度的窖藏，这一章的几条线索，在后面的章里还会出现。
\begin{itemize}
\item \textbf{金属是货物还是资产。}贵金属的流出，不等于经常账户的逆差（\chref{ch:external}）。罗马的钱币运出去，可以记作出口的货物，也可以记作减少的资产，式~\eqref{eq:india-coin}。今天的规则只把货币当局持有的储备黄金当作金融资产。
\item \textbf{风险与利率。}《政事论》按借款人分档的利率，用式~\eqref{eq:africa-risk} 读，是不同的还款概率；海事贷款把船沉的风险放进利率。地中海的商人后来用合伙来分担同样的风险（\chref{ch:islamic}）。
\item \textbf{印度洋的商路。}季风把红海、阿拉伯与印度连成一个一年一个来回的市场；伊斯兰世界兴起以后，把这条商路与地中海接在一起（\chref{ch:islamic}）。
\item \textbf{把金银流出当作损失。}普林尼的抱怨，一千多年后又出现在反对东印度公司运银出口的声音里，孟用贸易差额论回应它（\chref{ch:mercantilism}）。金银的流动怎样调整物价与贸易，见\chref{ch:bullion}。
\item \textbf{成色与信任。}印度人只要改革以前的旧银币。在法令管不到的地方，钱按所含的金属作价，外观只有在能说明成色时才值钱（\chref{ch:greekroman}的式~\eqref{eq:greek-value}）。今天的货币在国际上的地位，也要靠信任来维持（\chref{ch:hierarchy}）。
\end{itemize}

\Needspace{9\baselineskip}
\begin{takeaways}
\begin{itemize}
\item 季风使罗马与印度之间一年一个来回：按普林尼的日程，从亚历山大里亚到穆齐里斯 94 天，十二月到一月中旬回航；斯特拉波说约有一百二十条船从米奥斯·霍尔莫斯开往印度。汉朝的南海航路去程至少一年，靠“蛮夷贾船，转送致之”。
\item 《政事论》是一部年代不定的治国手册（Fleet：公元前 321—前 296 年；Olivelle：50—125 年成书、175—300 年改编）。它要国家专营矿与盐、集中商品定价、留一半储备防灾；向国库交钱要多付 13.125\%；它优待进口商，只问对外贸易“还有没有利润的余地”。
\item 它的利率按借款人分档，年率 15\%、60\%、120\%、240\%。按式~\eqref{eq:africa-risk}、以 15\% 为无风险利率，海上商人的还款概率只有 0.34；这么高的利率不只补偿船沉，也可能只是法定的上限。
\item 赫尔玛波隆号一条船的货值 692.7 万德拉克马（约同数的塞斯特斯），要交四分之一的关税 173.2 万。普林尼说印度每年吸走的钱，本章按拉丁原文读作五千万（英译作五亿五千万，与 12.84 对不上），三个地区合计一亿，只占罗马 GDP 的 0.3\%—0.6\%。
\item 钱币当作货物，用钱币买胡椒的经常账户是平衡的；当作资产，才有逆差，式~\eqref{eq:india-coin}。印度把罗马钱当作商品，要尼禄改革以前的旧银币（能认出发行者的 1 世纪银币 99.3\% 是奥古斯都与提比略的），后来要金币。出土的钱只能说明构成，不能量出流量。
\end{itemize}
\end{takeaways}

\begin{selfcheck}
\begin{enumerate}
\item 按普林尼的日程，从贝勒尼基到穆齐里斯要多少天？回程最晚在什么时候启航？
\item 《政事论》“林中”的利率折成年率是多少？以 15\% 为无风险利率，按式~\eqref{eq:africa-risk}，债主估计的还款概率是多少？
\item 如果赫尔玛波隆号只运了象牙一批（76 塔兰特 5675 德拉克马），交了四分之一的关税以后，货主还剩多少德拉克马？
\item 如果罗马的 GDP 是 200 亿塞斯特斯（Scheidel 说“略少于 200 亿”），普林尼的一亿占多少？印度的五千万呢？
\item 罗马用 100 万塞斯特斯的钱币买了胡椒，没有别的交易。按式~\eqref{eq:india-coin} 的两种记法，经常账户各是多少？按今天的规则应当用哪一种？
\item 为什么印度出土的 1 世纪银币几乎都是奥古斯都与提比略的？用式~\eqref{eq:greek-value} 解释。这个解释能说明全部吗？
\item 普林尼的“五亿五千万”与“一百倍”，哪一个可以通过重读原文改正，哪一个只能当作修辞？为什么？
\end{enumerate}
\end{selfcheck}

\begin{answers}
\item $30+40=70$ 天。回程最晚在埃及梅基尔月的第六天（罗马历一月的伊都斯日）以前启航。
\item $10\times12=120\%$；$p=1.15/2.20=0.52$。
\item $461675\times3/4=346256.25$ 德拉克马。
\item 一亿是 $1/200=0.5\%$；五千万是 0.25\%。
\item 当作货物：出口 100 万（钱币）、进口 100 万（胡椒），$\mathit{CA}=0$；当作罗马的金融资产：$\mathit{CA}=-100$ 万，用对外资产减少 100 万来弥补。罗马商人手里的钱币不是货币当局的储备，按今天的规则与非货币黄金一样记作货物，用第一种。
\item $v=c+\tau(\bar v-c)$。在印度，罗马的法令管不着，没有人保证 $\bar v$，钱按所含的银作价（$\tau=0$，$v=c$），但 $c$ 要知道。尼禄 64 年的改革使第纳尔的含银从约 90\% 降到 80\%，此后新旧两种标准并用，模样相近的钱成色不同，新钱要重新验；旧钱一看就知道成色，所以运到印度的是旧钱，罗马商人后来改用金币。这不是全部：同在改革以前的卡利古拉与克劳狄乌斯，出土的银币合起来也只有 14 枚。
\item “五亿五千万”可以改：拉丁原文是一个框起来的数码 D，按罗马人以十万为单位的计法是五千万（校勘记另有 DL，五千五百万）；英译多读了十倍，而且与 12.84 三个地区合计“按最低的估计”只有一亿对不上。“一百倍”改不了：原文就是 centiplicato，按它算，印度货在罗马的价值是 GDP 的 27.8\%，不可能，只能当作修辞。
\end{answers}

回到普林尼的那笔账。运往印度的钱币，是装在船上运过大洋的；据 Casson 的推断，赫尔玛波隆号的那船货，也是借钱买来、再押给债主的。这一章的商人用钱币付账，用货物作抵押，借着季风一年走一个来回，路上要防海盗，还要担心沉船。伊斯兰世界兴起以后，这条商路与地中海接在了一起。在那里做生意的商人，能不能不运金属，就把钱付到远方、为远航筹到资本？在禁止利息的规则下，他们又怎样放贷（\chref{ch:islamic}）？
